\documentclass[12 pt, letterpaper]{hmcpset}
\usepackage{hw-preamble}

\name{Jayden Thadani}
\class{MATH 176 --- Algebraic Geometry}
\assignment{Homework assignment \#10}
\duedate{15th November 2024}

\begin{document}

\begin{problem}[Problem 1.]
	Let $\mathcal{O}$ be an integral domain.

	\begin{enumerate}
		\item Consider a chain $I_{1} \subseteq I_{2} \subseteq \dots I_{k} \dots \subsetneq \mathcal{O}$ of proper ideals of $\mathcal{O}$. Show that $I = \bigcup_{k \ge 1} I_{k}$ is also a proper ideal of $\mathcal{O}$.
		\item Show that each proper ideal $I$ of $\mathcal{O}$ is contained in a maximal ideal $M$ of $\mathcal{O}$.
	\end{enumerate} 
\end{problem}

\begin{solution}
	\begin{enumerate}
		\item To show $I$ is an ideal, we need to show that for any $i_{1}, i_{2} \in I$, and any $r \in \mathcal{O}$, we have $i_{1} + i_{2}, r i_{1} \in I$. If $i_{1}, i_{2} \in I$, then they must be in some $I_{k_{1}}$ and $I_{k_{2}}$ respectively. If we assume, without loss of generality that $k_{1} \le k_{2}$, then $I_{k_{1}} \subseteq I_{k_{2}}$ so $i_{1}, i_{2} \in I_{k_{2}}$ giving us $i_{1} + i_{2} \in I_{k_{2}} \subseteq I$. Further, $r i_{1} \in I_{k_{1}} \subseteq I$. Thus, $I$ is an ideal of $\mathcal{O}$.

			Note that $1 \notin I_{k}$ for all $k$ since they are all proper ideals, so $1 \notin I$. Thus, $I$ is a proper ideal.
		\item Part (a) is exactly the hypothesis of Zorn's lemma --- every chain of proper ideals, is contained in the upper bound proper ideal that is the union of the ideals. Thus, every ideal is contained in a maximal proper ideal.
	\end{enumerate}
\end{solution}

\pagebreak

\begin{problem}[Problem 2.]
	For a prime ideal $I$ of an integral domain $\mathcal{O}$, define $\operatorname{ht} I$, the \emph{height} of $I$ as the largest $m$ such that there is a chain $I_{0} \subsetneq I_{1} \subsetneq \dots \subsetneq I_{m}$ of primes $I_{k}$ satisfying $I_{0} = \{ 0 \}$ and $I_{m} = I$.
	\begin{enumerate}
		\item Given a prime ideal $I$, show there exists a maximal ideal $M$ such that $\operatorname{ht} M \ge \operatorname{ht} I$.
		\item Show that $\operatorname{ht} I = 0$ if and only if $I$ is the zero ideal.
		\item Assume that $\mathcal{O}$ is a unique factorisation domain. Show that a non-zero prime has $\operatorname{ht} 1$ if and only if $I = \left < f \right >$ for some irreducible element $f \in \mathcal{O}$.
	\end{enumerate}
\end{problem}

\begin{solution}
	\begin{enumerate}
		\item Every prime ideal $I$ is contained in a maximal ideal $M$. Thus, there's a chain
			\[
				I_{0} \subsetneq I_{1} \subsetneq \dots \subsetneq I_{\operatorname{ht} I} \subsetneq M.
			\]
			Thus, $\operatorname{ht} M \ge \operatorname{ht} I + 1 \ge \operatorname{ht} I$.
		\item If $I \neq \{ 0 \}$, then there exists a chain of length $1$ ---
			\[
				I_{0} \subsetneq I	
			\]
			so $\operatorname{ht} I \ge 1$ and is not $0$.
		\item Suppose $\operatorname{ht} I = 1$. By unique factorisation, there exists an irreducible $f \in I$ (just factorise any $i \in I$ into irreducible elements and apply the primeness of $I$). Clearly
			\[
				I_{0} \subsetneq \left < f \right > \subseteq I.
			\]
			Since $\operatorname{ht} I = 1$, this chain must have length $1$ with $I = \left < f \right >$.

			Suppose $I = \left < f \right >$ for an irreducible $f$. If we have a chain
			\[
				I_{0} \subsetneq I_{1} \subseteq \dots \subseteq I_{m} = I
			\]
			then by unique factorisation there is an irreducible $f_{1} \in I_{1}$ giving us the chain
			\[
				I_{0} \subsetneq \left < f_{1} \right > \subseteq I_{1} \subseteq \dots \subseteq I_{m} = \left < f \right >.
			\]
			Thus, $f \mid f_{1}$. But $f_{1}$ and $f$ are irreducible, so $f_{1} = u f$ for some unit $u$. Thus, $\left < f \right > = \left < f_{1} \right >$ and the chain reduces to
			\[
				I_{0} \subsetneq \left < f_{1} \right > = I_{1} = \dots = I_{m} = \left < f \right >.
			\]
			That is, any chain of length $n$ reduces to a proper chain of length $1$, so $\operatorname{ht} I = 1$.
	\end{enumerate}
\end{solution}

\pagebreak

\begin{problem}[Problem 3.]
	Let $\mathcal{O}$ be an integral domain. Define $\operatorname{dim} \mathcal{O}$, the \emph{Krull dimension} of $\mathcal{O}$ as the largest $\operatorname{ht} I$ ranging over $I \in \operatorname{Spec} \mathcal{O}$.
	\begin{enumerate}
		\item Show that $\dim \mathcal{O} = 0$ if and only if $\mathcal{O}$ is a field.
		\item Show that the following are equivalent:
			\begin{enumerate}[label = (\roman*)]
				\item $\dim \mathcal{O} = 1$.
				\item Every non-zero $I \in \operatorname{Spec} \mathcal{O}$ has $\operatorname{ht} I = 1$.
				\item Every prime ideal $I$ of $\mathcal{O}$ is either the zero ideal or a maximal ideal.
			\end{enumerate}
	\end{enumerate}
\end{problem}

\begin{solution}
	\begin{enumerate}
		\item $\mathcal{O}$ is a field if and only if its only ideals are $\mathcal{O}$ and $\{ 0 \}$ --- since $\mathcal{O} / \{ 0 \} \cong \mathcal{O}$ is a field when $\mathcal{O}$ is a field, the zero ideal must be maximal in any field.

			Thus, if $\mathcal{O}$ is a field, the only prime ideal is $\{ 0 \}$ and the largest height of a prime ideal is $0$.

			If the largest $\operatorname{ht} I$ for  $I \in \operatorname{Spec} \mathcal{O}$ is $0$, then the largest maximal ideal height is also $0$. That is, $\{ 0 \}$ is the only maximal ideal. Thus, $\mathcal{O} \neq \mathcal{O} / \{ 0 \}$ is a field.
		\item We can see that (i) and (ii) are equivalent since $\dim \mathcal{O} = 1$ is equivalent to $\operatorname{ht} I \le 1$ and the existence of a non-zero prime $I$ (for which we already have $0 < \operatorname{ht} I$).

			Now we show that (ii) and (iii) are equivalent. Suppose every non-zero prime $I$ has height $1$. The maximal ideal $M$ containing $I$ is also a prime, and so has height $1$. Since
			\[
				I_{0} \subsetneq I \subseteq M
			\]
			we must have $I = M$. That is, $I$ is its own maximal ideal --- $I$ is maximal.

			Suppose every non-zero prime ideal is maximal. Then for any chain of primes
			\[
				I_{0} \subseteq I_{1} \subseteq \dots \subseteq I_{m}
			\]
			since $I_{1}$ is maximal, $I_{1} = \dots = I_{m}$. Thus, the longest proper chain of primes has length at most $1$.
	\end{enumerate}
\end{solution}

\pagebreak

\begin{problem}[Problem 4.]
	Let $R = F[x_{1}, x_{2}, \dots, x_{n}]$ be the polynomial ring in $n$ variables over an algebraically closed $F$.
	\begin{enumerate}
		\item Show that $\dim \mathcal{O} = \dim R - \operatorname{ht} I$ whenever $\mathcal{O} = R / I$ for some prime ideal $I$.
		\item Show $\dim R = n$.
		\item Show $\operatorname{ht} M = n$ for any $M \in \operatorname{mSpec} R$.
	\end{enumerate}
\end{problem}

\begin{solution}
	\begin{enumerate}
		\item By the fourth isomorphism theorem, we have an inclusion preserving bijection between the ideals of $\mathcal{O}$ and the ideals of $R$ that contain $I$. Further, using the third isomorphism theorem we can show that this bijection preserves primeness. Specifically, for any $P$ with $I \subseteq P$ we have $\frac{\mathcal{O}}{P / I} \cong R / P$. Now since $\frac{\mathcal{O}}{P / I}$ is an integral domain if and only if $R / P$ is an integral domain, $P / I$ is a prime ideal if and only if $P$ is a prime ideal.

			Thus, for any proper chain of prime ideals
			\[
				\{ 0 \} = \overline{J}_{0} \subsetneq \overline{J}_{1} \subsetneq \dots \subsetneq \overline{J}_{k} = \overline{J}
			\]
			we have a chain of prime ideals
			\[
				I = J_{0} \subsetneq J_{1} \subsetneq \dots \subsetneq J_{k} = J
			\]
			where $J_{i} / I = \overline{J}_{i}$. Since there exists a chain of prime ideals up to $I$ with maximal length $\operatorname{ht} I$, and the maximum length of a chain of prime ideals in $R$ is $\dim R$, the maximal length of a chain of ideals in $\mathcal{O}$ must be the difference. That is, $\dim \mathcal{O} = \dim R - \dim I$.
		\item We show that $\dim F[x_{1}, \dots , x_{n}] = n$ by induction. Note that we have $\dim F = 0$ since the only prime ideal of $F$ is the zero ideal. Suppose $\dim F[x_{1}, \dots, x_{n - 1}] = n - 1$. Notice that $\left < x_{n} \right >$ is prime in $F[x_{1}, \dots, x_{n}]$ with height $1$ by problem 2. Since
			 \[
				F[x_{1}, \dots, x_{n}] / \left < x_{n} \right > \cong F[x_{1}, \dots, x_{n}]
			\]
			we have
			\[
				\dim F[x_{1}, \dots, x_{n}] - \operatorname{ht} \left < x_{n} \right > = \dim F[x_{1}, \dots, x_{n}].
			\]
			Rearranging, we get
			\[
				\dim F[x_{1}, \dots, x_{n}] = n - 1 + 1 = n.	
			\]
		\item We know that the maximal ideals of $R$ are all of the form
			\[
				M = \left < x_{1} - a_{1}, \dots, x_{n} - a_{n} \right >
			\]
			for $a_{j} \in F$. Consider the chain
			\[
				\{ 0 \} = I_{0} \subsetneq \left < x_{1} - a_{1} \right > \subsetneq \left < x_{1} - a_{1}, x_{2} - a_{2} \right > \subsetneq \dots \subsetneq M.
			\]
			Since each $I_{i} = \left < x_{1} - a_{1}, \dots, x_{i} - a_{i} \right >$ quotients out to give an integral domain $R / I_{k} \cong F[x_{i + 1}, \dots, x_{n}]$, each $I_{k}$ is prime. Thus, $\operatorname{ht} M \ge n$. But since $\dim R = n$, we also have $\operatorname{ht} M \le n$. Thus, $p\operatorname{ht} M = n$.
	\end{enumerate}
\end{solution}

\end{document}
